\chapter*{Abstract}
This project develops a digital ecosystem for glycemic monitoring and therapeutic education by integrating Near Field Communication (NFC) technology and Clinical Decision Support Systems (CDSS). The system enables real-time interstitial glucose data acquisition by connecting directly to continuous monitoring sensors (FreeStyle Libre 2 Plus), centralizing critical records such as insulin profiles, medical guidelines, carbohydrate intake, and physical activity.

The innovative core of the project lies in a conversational interface based on Generative Artificial Intelligence (GPT-style model). This module functions as a pedagogical assistant that analyzes real-time variables, such as Insulin on Board (IOB), to provide personalized guidance in risk situations, including the prevention of post-exercise hypoglycemia or the management of basal insulin.

With a strictly educational focus, the primary objective is to empower users without prior knowledge, allowing them to consolidate concepts and learn intuitively through their own data. The system prioritizes safety by contextualizing every AI response as complementary support that reinforces, rather than replaces, established medical guidelines. Ultimately, this project demonstrates how the convergence of biometric monitoring and AI facilitates a deeper and more accessible understanding of daily diabetes management.
\vspace{.5cm}

\textbf{Keywords:} Diabetes Mellitus, Generative Artificial Intelligence, Continuous Glucose Monitoring (CGM), NFC, Clinical Decision Support Systems (CDSS), Patient Education, mHealth.